% =====================================================================
\chapter{Marco teórico}\label{cap:marco-teorico}
% Plantilla ITBA: sustento formal que valida las decisiones de diseño;
% evitar el "efecto enciclopedia".
% =====================================================================

% ---------------------------------------------------------------------
\section{Ventanas críticas del desarrollo hipocampal: fundamento del diseño etario y del período de lavado}\label{sec:ventanas}

La elección de dos ventanas de exposición cualitativamente distintas ---PD25 y PD75--- no es arbitraria, sino que se apoya en la trayectoria de maduración del hipocampo. La formación hipocampal, y en particular el giro dentado, presenta un desarrollo postnatal prolongado: la densidad de espinas dendríticas en neuronas CA1 se incrementa de forma pronunciada entre PD7 y PD28, la inducibilidad de la potenciación a largo plazo alcanza su máximo entre PD10 y PD60, y la neurogénesis granular es máxima en el período postnatal temprano \cite{sakimoto2022critical}. PD25 se ubica, por tanto, dentro de una ventana de alta plasticidad en la que el circuito aún se está cableando y resulta máximamente susceptible a la programación por factores ambientales; PD75, en cambio, corresponde a un circuito estructuralmente maduro. Contrastar ambos puntos de inicio operacionaliza la hipótesis de los orígenes del desarrollo de la salud y la enfermedad (DOHaD), según la cual una misma exposición produce consecuencias divergentes en función del estadio madurativo en que incide \cite{bhatt2025prenatal,coirini2022sucrose}.

El período de lavado ---acceso exclusivo a agua tras la exposición, hasta PD75 en el grupo juvenil y PD125 en el adulto--- constituye el elemento de diseño que confiere validez a la inferencia. Durante la exposición coexisten efectos de dos naturalezas: perturbaciones agudas de origen metabólico o farmacológico, presentes solo mientras dura el consumo, y una programación persistente de tipo estructural, que sobrevive a la retirada del estímulo. Al interponer el lavado entre el fin del tratamiento y la evaluación conductual e histológica, el diseño aísla el segundo componente ---el fenómeno de interés, es decir, la modificación duradera del sustrato neural--- del primero, que actuaría como confusor. Sin esta separación temporal, cualquier diferencia observada sería ambigua respecto de su carácter transitorio o permanente.

% ---------------------------------------------------------------------
\section{Bases farmacológicas y metabólicas del esquema de exposición}\label{sec:bases}
% Nota de trabajo: queda abierto un problema de encuadre metodológico en esta sección.

El esquema de tres brazos ---stevia, sacarosa y agua--- responde a una lógica de anclaje experimental. La sacarosa al 10\,\% p/v reproduce el patrón de consumo de bebidas azucaradas y posee un fenotipo adverso documentado sobre memoria, conducta ansiosa y neurogénesis \cite{reichelt2016sucrose,coirini2022sucrose,kruse2019sucrose}; funciona, por tanto, como control positivo, esto es, como referencia de un efecto perjudicial conocido contra el cual calibrar la respuesta a la stevia. El agua provee la línea de base. La stevia al 4\,\% p/v, preparada como decocción de hojas según el método de Sharma et~al.\ \cite{sharma2009stevia}, se ajusta para igualar la intensidad edulcorante de la sacarosa al 10\,\% p/v sin aportar calorías \confirmar{cita que respalde la equivalencia de dulzor entre la stevia al 4\,\% p/v y la sacarosa al 10\,\% p/v}. Esta equivalencia sensorial con aporte calórico nulo es precisamente la manipulación experimental crítica: permite disociar la vía del sabor dulce y la señalización central de la vía calórica y metabólica propia de la sacarosa. En consecuencia, comparar la stevia simultáneamente contra ambos anclajes ubica su efecto sobre un eje que va de lo inocuo a lo perjudicial, lo que ningún diseño de dos brazos podría resolver.

La farmacocinética de los glucósidos de esteviol justifica, además, por qué el carácter no calórico no equivale a inocuidad. Estos compuestos no se absorben intactos en el tracto superior; son hidrolizados por la microbiota colónica a esteviol, con la consiguiente posibilidad de efectos mediados por el eje intestino-cerebro y por la modulación de la homeostasis glucídica \cite{delagarza2022maternal,tsan2022lcs,noble2021microbial}. A ello se suma la evidencia de acciones insulino-miméticas y antioxidantes directas de los glucósidos de esteviol sobre tejidos periféricos \cite{prata2017glycosides}. Este perfil ---no calórico pero biológicamente activo--- es el que hace necesaria una caracterización sistemática y no una asunción a priori de neutralidad.

% ---------------------------------------------------------------------
\section{La cascada neurogénica y la separación de patrones: fundamento de los marcadores y de la batería conductual}\label{sec:cascada}

La selección conjunta de tres marcadores inmunocitoquímicos ---PCNA, DCX y NeuN--- se fundamenta en que la neurogénesis del giro dentado es un proceso estadificado, no un evento único. Los progenitores atraviesan una progresión ordenada: células madre y progenitores en proliferación (PCNA$^{+}$), neuroblastos y neuronas inmaduras en migración y diferenciación (DCX$^{+}$) y neuronas granulares maduras post-mitóticas integradas al circuito (NeuN$^{+}$). Cada marcador indexa, así, una etapa distinta de la cascada. Emplear los tres ---y no uno solo--- permite localizar en qué punto del proceso actúa una eventual perturbación (por ejemplo, una reducción de la proliferación frente a un fallo en la supervivencia o la maduración), resolución que un marcador aislado no puede proporcionar. La lectura histológica adquiere de este modo un carácter cuantitativo y por etapas, condición necesaria para vincular la estructura con la función conductual.

\begin{figure}[htbp]
  \centering
  \incluirfigura[width=0.85\linewidth]{Fig3_cascada_neurogenica.png}
  \caption[Esquema de la cascada neurogénica del giro dentado]{Esquema de la cascada neurogénica del giro dentado y sus marcadores. \pendiente{Insertar esquema: diagrama de flujo simple con tres o cuatro etapas en secuencia ---célula madre/progenitor en proliferación $\to$ neuroblasto/neurona inmadura en migración $\to$ neurona granular madura integrada al circuito--- etiquetando cada etapa con su marcador correspondiente (PCNA$^{+}$, DCX$^{+}$, NeuN$^{+}$). Estilo minimalista, coherente con la Figura~\ref{fig:linea-temporal}.}}
  \label{fig:cascada-neurogenica}
\end{figure}

Ese vínculo se articula formalmente a través del concepto de separación de patrones. Desde una perspectiva computacional, la separación de patrones es la transformación de representaciones de entrada solapadas en representaciones de salida más ortogonales, lo que reduce la interferencia entre memorias similares; el giro dentado se postula como el sustrato de esta operación mediante un recodificado disperso y de alta dimensionalidad, en el que las neuronas granulares de nueva generación resultarían críticas para discriminar contextos espaciales altamente superpuestos \cite{reichelt2016sucrose,reichelt2021slr}. La prueba SLR operacionaliza conductualmente esta función a través de dos configuraciones que gradúan la demanda de separación: la configuración d-SLR, con objetos dispuestos de forma equidistante (bajo solapamiento espacial), impone una demanda baja, mientras que la s-SLR, con objetos próximos entre sí (alto solapamiento), impone una demanda elevada que tributa selectivamente sobre la separación dependiente de neurogénesis \cite{reichelt2021slr}. La s-SLR constituye, por tanto, la sonda sensible que conecta el desempeño conductual con la lectura histológica de DCX y PCNA.

\begin{figure}[htbp]
  \centering
  \incluirfigura[width=0.95\linewidth]{Fig3_SLR_carga_cognitiva.png}
  \caption[Gradiente de carga sobre la separación de patrones en la prueba SLR]{Gradiente de carga sobre la separación de patrones en la prueba SLR, según la proximidad de los objetos durante la fase de muestreo. A menor distancia entre dos de los objetos, mayor es la demanda de plasticidad hipocampal necesaria para distinguir sus ubicaciones. Este trabajo emplea únicamente las configuraciones d-SLR y s-SLR; la configuración xs-SLR (carga \enquote{muy alta}) se muestra por completitud del esquema original pero no fue evaluada. Adaptado de \cite{ghoshswaby2021slr}.}
  \label{fig:slr-carga}
\end{figure}

El Open Field y el NOR completan la batería probando construcciones disociables ---distribución espacial y locomoción como índices de conducta tipo ansiosa el primero; preferencia por la novedad con participación perirrinal e hipocampal el segundo \cite{leger2013nor}---, de modo que el conjunto cubre dominios conductuales distintos y no redundantes. Las tasas de exploración del NOR y la SLR se definen como el cociente $t_N/(t_N+t_F)$, acotado en $[0,1]$ y con 0,5 como valor de azar, lo que provee un observable normalizado y adimensional, comparable entre animales.

% ---------------------------------------------------------------------
\section{Fundamento analítico del abordaje multivariado}\label{sec:multivariado}

Las métricas que \anymaze{} extrae de cada prueba ---distancia recorrida, velocidad media, tiempo y distancia por zona, tiempo explorando cada objeto, cantidad de veces de exploración de objetos--- constituyen un vector de variables fuertemente correlacionadas. Analizarlas mediante ANOVA univariado, una por una, presenta limitaciones formales: la realización de $m$ contrastes infla la tasa de error por familia,
\begin{equation}
  \mathrm{FWER} = 1 - (1-\alpha)^{m},
  \label{eq:fwer}
\end{equation}
y su corrección por métodos tipo Bonferroni resulta excesivamente conservadora justamente cuando las variables están correlacionadas; el análisis variable por variable descarta la estructura de covarianza conjunta, de modo que una diferencia entre grupos puede residir en una combinación lineal de variables sin que ninguna alcance significancia de forma aislada; y los supuestos de normalidad y homocedasticidad son frecuentemente violados por datos conductuales de conteo y tiempo. El abordaje multivariado no es, por tanto, un adorno metodológico, sino la respuesta pertinente a la estructura de los datos.

El análisis de componentes principales (PCA) aborda la redundancia \cite{jolliffe2016pca}. Estandarizadas las variables, se resuelve la descomposición espectral de la matriz de correlación,
\begin{equation}
  \mathbf{R} = \mathbf{V}\boldsymbol{\Lambda}\mathbf{V}^{\mathsf{T}},
  \label{eq:pca}
\end{equation}
cuyos autovectores definen las direcciones principales y cuyos autovalores $\lambda_k$ cuantifican la varianza explicada; la proporción $\lambda_k/\sum_j \lambda_j$ y la varianza acumulada guían la retención de las primeras componentes. Como las métricas de \anymaze{} son colineales, el PCA concentra la señal en pocas dimensiones ortogonales, mitigando la redundancia que penaliza al contraste univariado. Sobre las puntuaciones resultantes se aplica clustering jerárquico con criterio de Ward, que minimiza el incremento de la varianza intra-grupo al fusionar clusters \cite{ward1963hierarchical}; el dendrograma permite examinar si los perfiles conductuales segregan en grupos naturales sin imponer a priori las etiquetas de exposición. Descubrir esa estructura sin supervisión es lo que da sustento al objetivo de máxima M1: que la separación entre tratamientos emerja de los propios datos constituye una afirmación más fuerte que una comparación entre grupos previamente rotulados.

El objetivo M2 ---determinar si las métricas conductuales predicen el grupo de exposición y con qué peso relativo--- es un problema de clasificación supervisada con estimación de importancia de variables, para el cual el Random Forest resulta el estimador adecuado \cite{breiman2001rf}. El modelo agrega árboles de decisión entrenados sobre muestras \textit{bootstrap} y con selección aleatoria de variables en cada nodo, lo que decorrelaciona los árboles y controla la varianza; las observaciones no incluidas en cada \textit{bootstrap} (\textit{out-of-bag}) proveen una estimación insesgada del error de generalización sin requerir un conjunto de prueba independiente \cite{breiman2001rf}, ventaja decisiva ante los tamaños muestrales acotados de la experimentación animal \cite{charan2013sample}. La importancia relativa de cada métrica se cuantifica por la disminución media de impureza o por la caída de exactitud al permutar aleatoriamente sus valores, lo que responde directamente a M2, y el modelo acomoda relaciones no lineales, interacciones y colinealidad sin exigir supuestos distribucionales. En síntesis, PCA, clustering jerárquico y Random Forest no compiten con el análisis estadístico clásico ---que conserva su lugar para los contrastes univariados planificados--- sino que lo complementan allí donde la naturaleza multivariada y correlacionada de los datos conductuales excede su alcance; esta articulación constituye el aporte diferencial del proyecto desde la bioingeniería.
